%% Fulcrum — Optimal Per-Region DP Noise Allocation for Hierarchical Federated Learning
%% Build: latexmk -pdf main.tex

\documentclass[letterpaper,journal]{IEEEtran}

% --- Math ---
\usepackage{amsmath,amsfonts,amssymb}
\usepackage{amsthm}

% --- Algorithms ---
\usepackage{algorithmic}
\usepackage{algorithm}

% --- Tables and figures ---
\usepackage{array}
\usepackage[caption=false,font=normalsize,labelfont=sf,textfont=sf]{subfig}
\usepackage{booktabs}
\usepackage{multirow}
\usepackage{graphicx}
\usepackage{stfloats}
\usepackage{threeparttable}

% --- Other ---
\usepackage{textcomp}
\usepackage{url}
\usepackage{xcolor}
\usepackage{orcidlink}
\usepackage{xspace}

% --- Theorem environments ---
\newtheorem{theorem}{Theorem}[section]
\newtheorem{lemma}[theorem]{Lemma}
\newtheorem{corollary}[theorem]{Corollary}
\newtheorem{proposition}[theorem]{Proposition}
\newtheorem{definition}[theorem]{Definition}

\hyphenation{op-tical net-works semi-conduc-tor}

% --- Macros ---
\newcommand{\eg}{\textit{e.g.,}\xspace}
\newcommand{\ie}{\textit{i.e.,}\xspace}
\newcommand{\Fulcrum}{\textsc{Fulcrum}\xspace}
\newcommand{\KL}{D_{\mathrm{KL}}}

\begin{document}

\title{Optimal Noise Allocation for Hierarchical\\ Federated Learning}

\author{Murtaza Rangwala\,\orcidlink{0000-0000-0000-0000},~\IEEEmembership{Member,~IEEE,}
        Richard O. Sinnott\,\orcidlink{0000-0000-0000-0000},
        and~Rajkumar Buyya\,\orcidlink{0000-0000-0000-0000},~\IEEEmembership{Fellow,~IEEE}%
\thanks{The authors are with the Quantum Cloud Computing and Distributed Systems (qCLOUDS) Lab, School of Computing and Information Systems, The University of Melbourne, Australia.}}

\markboth{}{}
\maketitle

\begin{abstract}
Hierarchical federated learning inserts regional aggregators between clients and
the cloud, so that a participant's update is observed only after it has been
combined with those of its neighbours. The size of that aggregation region
determines how much protection the region itself provides, and regions in real
deployments are rarely the same size. Standard practice nonetheless applies a
single noise multiplier to every participant, which necessarily calibrates for
the most exposed region and over-provisions all the others. We show that the
resulting allocation problem has a closed-form solution. We first derive a
per-client bound on the mutual information between a participant's local class
distribution and any estimate an observer can form of it, taking care that the
adjacency notion in the mechanism matches the quantity being protected. We then
solve for the noise allocation that minimises the worst-case value of this bound
under a fixed utility budget, and prove it is min-max optimal. The saving over
uniform allocation is exactly $\delta = 1 - \langle \rho \rangle_V / \rho_{\max}$,
where $\rho_r$ is the inverse participation ratio of region $r$. This quantity
depends only on the deployment's region-size profile, so it can be computed
before any training takes place, and it is zero precisely when all regions have
the same effective size. Experiments on two modalities, image classification with
a frozen ResNet-18 and text classification with a frozen sentence encoder, confirm
the prediction at $\varepsilon = 0.99$. Accuracy at matched per-client privacy
improves by up to $13.75$ percentage points where $\delta$ is large, is exactly
zero in two balanced control deployments where the theory predicts no gain, and is
negative for a control that uses the same noise dispersion but assigns it to the
wrong regions.
\end{abstract}

\begin{IEEEkeywords}
Federated learning, differential privacy, hierarchical aggregation, noise allocation.
\end{IEEEkeywords}

\section{Introduction}
\label{sec:intro}

\IEEEPARstart{F}{ederated} learning allows several organisations to train a
shared model without pooling their data. In cross-silo and edge deployments the
participants seldom report to a single server. Instead they report to an
intermediate aggregator, which combines the updates it receives and forwards only
the combination. Cellular deployments place these aggregators at base stations,
clinical consortia place them at the hospital or country level, and industrial
deployments place them at regional gateways. This three-tier arrangement is
usually adopted for bandwidth and scalability reasons, but it also changes the
privacy picture, because an observer no longer sees any single participant's
update in isolation.

The size of the group a participant is aggregated with therefore matters. A silo
in a region of twenty peers is concealed by the combined noise of twenty
participants, while a silo that shares its region with only one other is
concealed by far less. Real deployments are heterogeneous in exactly this
respect. Base station loads are long tailed, with dense urban cells serving many
devices and rural cells serving very few. Multinational consortia contribute
wildly uneven numbers of sites per country.

Standard practice ignores this. Differentially private federated learning applies
one noise multiplier to every participant, which must be calibrated so that the
most exposed participant is adequately protected. Every participant in a larger
region then receives more noise than its own privacy requirement calls for, and
that surplus noise is paid for in model accuracy by the whole federation.

This paper asks how noise should be distributed across regions instead, and shows
that the question has a clean answer. Our contributions are as follows.

We derive a per-client bound on the information an observer can obtain about a
participant's local class distribution (Theorem~\ref{thm:bound}). The bound is
stated for silo-level adjacency, meaning that the mechanism protects a
participant's entire dataset rather than any single record in it. This matters
because the quantity we wish to protect, the proportion of a silo's data
belonging to a sensitive class, is a property of the dataset as a whole and is not
controlled by record-level guarantees. The bound also carries an explicit cap at
the entropy of the target, which makes it possible to check whether the guarantee
is informative at a given operating point.

We then pose the allocation problem, minimising the worst-case value of that
bound over all participants subject to a fixed budget on the noise injected into
the global model, and give its closed-form solution
(Theorem~\ref{thm:allocation}). The correct budget is the noise variance actually
entering the aggregate, which under weighted averaging is
$\sum_i w_i^2 \sigma_i^2$ rather than $\sum_i \sigma_i^2$.

The improvement over uniform allocation follows in closed form
(Corollary~\ref{cor:delta}). Writing $\rho_r$ for the inverse participation ratio
of region $r$, the fraction of the noise budget saved at equal per-client privacy
is $\delta = 1 - \langle \rho \rangle_V / \rho_{\max}$. Because $\delta$ is
determined by the region-size profile alone, a deployer can compute it in advance
and decide whether the method is worth applying. It is zero exactly when all
regions have the same effective size, which gives the approach an explicit
applicability test rather than a failure mode discovered after deployment.

Finally we validate the prediction experimentally on two modalities and report
the privacy parameters achieved in every run. The evaluation includes two
deployments for which the theory predicts no benefit at all, and a control that
uses the same noise dispersion but assigns it to the wrong regions. Both behave
as predicted, which we regard as more informative than the positive result alone.

\section{Related Work}
\label{sec:related}

\subsection{Differentially private federated learning}

Differentially private stochastic gradient descent~\cite{abadi2016dpsgd} clips
per-example gradients and adds calibrated noise, and its federated instantiation
adds noise to client updates before aggregation~\cite{mcmahan2018dprnn}. Privacy
accounting across rounds is usually performed with Renyi differential
privacy~\cite{mironov2017rdp} or the moments accountant. These mechanisms
distribute noise uniformly across participants.

\subsection{Heterogeneous and personalised privacy budgets}

A separate line of work assigns different budgets to different participants.
Jorgensen et al.~\cite{jorgensen2015personalized} introduced personalised
differential privacy, in which each individual declares a budget, and subsequent
work extended the idea to federated training. In these schemes the per-client
budget is an input to the system, supplied by the user. In ours it is an output
of an optimisation, determined by the participant's structural position. The two
are complementary, and an allocation combining declared preferences with
structural exposure is a natural extension that we do not pursue here.

\subsection{Clustered and hierarchical private federated learning}

Closest to our setting is work that combines hierarchical or clustered training
with differential privacy. Chandrasekaran et al.~\cite{chandrasekaran2022hfl}
propose hierarchical differential privacy, in which clients are organised into
zones with elected super-nodes and calibrated noise is added at the zone level.
Their per-zone noise scale depends on the number of participating clients in that
zone, and they observe that intermediate aggregation provides additional
protection at the central aggregator. We take the existence of this effect as our
starting point rather than as a contribution. What we add is the optimisation: we
show that choosing the per-region noise levels is a constrained min-max problem
with a closed-form solution, and we characterise exactly how much is gained
relative to uniform allocation.

The two settings also differ in where the noise originates, and the difference is
quantitative. If a trusted super-node adds a single noise draw on behalf of its
zone, the noise required per client scales as $1/m$ in a region of $m$ equally
weighted participants. If each participant adds its own noise before transmitting,
as we assume, the region's noise accumulates across $m$ independent draws and the
requirement scales as $1/\sqrt{m}$. The gap of $\sqrt{m}$ is the price of not
having to trust the aggregator, and we make the assumption explicit rather than
leaving it implicit in the mechanism.

Guo et al.~\cite{guo2025fedcdp} study clustered federated learning with
heterogeneous differential privacy. Their clusters group clients with similar data
distributions in order to train personalised models, and their intra-cluster
weighting adjusts aggregation weights as a function of a client's noise level
using a chosen logarithmic function. The direction is the reverse of ours: they
take the noise as given and choose the weights, whereas we take the structure as
given and choose the noise. Yin et al.~\cite{yin2026hdpfedcd} allocate noise
according to a per-example data-quality metric within each client's local dataset;
their topology is flat and their hierarchy is internal to a client, so region size
plays no role. Neither work states an optimality result or bounds the achievable
gain.

\subsection{Amplification by aggregation and by topology}

Secure aggregation~\cite{bonawitz2017secagg} conceals individual updates from the
server, and shuffling arguments quantify the resulting amplification, which
typically improves with the number of participants combined. Cyffers and
Bellet~\cite{cyffers2022amplification} and Cyffers et
al.~\cite{cyffers2022muffliato} develop network differential privacy for
decentralised averaging, where the guarantee between two nodes strengthens with
their distance in the communication graph. All of this work computes the
amplification enjoyed by a fixed, uniform mechanism. We treat the amplification
profile as the object to allocate against. When regions differ in effective size
the amplification is unequal, and uniform noise is then provably wasteful by a
quantity we compute exactly.

Table~\ref{tab:related} summarises the comparison.

\begin{table}[t]
\caption{Placement relative to the closest lines of work.}
\label{tab:related}
\centering
\small
\begin{tabular}{@{}lccc@{}}
\toprule
& Per-client & Allocation & Optimality \\
& noise varies & derived from & result \\
\midrule
DP-FedAvg~\cite{mcmahan2018dprnn}      & no  & ---                & --- \\
Personalised DP~\cite{jorgensen2015personalized} & yes & user preference & no \\
FedCDP~\cite{guo2025fedcdp}            & yes & noise level (heur.) & no \\
HDP-FedCD~\cite{yin2026hdpfedcd}       & yes & data quality (heur.) & no \\
Hierarchical DP~\cite{chandrasekaran2022hfl} & yes & zone size & no \\
Muffliato~\cite{cyffers2022muffliato}  & no  & ---                & --- \\
\textbf{This work}                     & \textbf{yes} & \textbf{zone exposure} & \textbf{yes} \\
\bottomrule
\end{tabular}
\end{table}

\section{System and Threat Model}
\label{sec:model}

\subsection{Federation}

A federation consists of $n$ silos $P_1, \dots, P_n$. Silo $i$ holds a local
dataset $D_i$ and contributes to the global model with aggregation weight
$w_i = |D_i| / \sum_j |D_j|$, so that $\sum_i w_i = 1$. Silos are partitioned
into disjoint aggregation regions. We write $A(i)$ for the region containing silo
$i$, and index regions by $r$. Training proceeds for $T$ rounds of federated
averaging.

The region structure is a property of the deployment. It may follow the physical
network, as when clients report to the base station serving them, or the
organisational structure, as when hospital sites report through a national
coordinator. We take it as given.

\subsection{Mechanism}

At round $t$, silo $i$ trains locally on $D_i$ starting from the current global
model $\theta^{(t-1)}$, forms its update $\Delta_i^{(t)}$, clips it so that
$\|\Delta_i^{(t)}\| \le C$, and adds Gaussian noise:
\begin{equation}
\tilde{\Delta}_i^{(t)} \;=\; \mathrm{clip}\big(\Delta_i^{(t)}, C\big)
\;+\; \xi_i^{(t)}, \qquad
\xi_i^{(t)} \sim \mathcal{N}\!\left(0, \sigma_i^2 C^2 I\right).
\label{eq:mechanism}
\end{equation}
The regional aggregator sums the noised updates it receives, and the observer
sees only the regional sums
\begin{equation}
Y_r^{(t)} \;=\; \sum_{j \in r} w_j \, \tilde{\Delta}_j^{(t)} .
\label{eq:observation}
\end{equation}

Two features of \eqref{eq:mechanism} deserve comment. First, the clipping is
applied to the silo's whole update rather than to per-example gradients, so the
mechanism protects the participation of an entire silo. This matches the quantity
we set out to protect, which is a property of the silo's dataset as a whole.
Second, each silo adds its own noise before transmitting, so the aggregator need
not be trusted to perform the noise addition honestly.

A useful practical consequence is that the amount of local computation performed
between communications does not affect the privacy accounting. Each silo releases
exactly one clipped and noised update per round regardless of how many local
optimisation steps produced it, so the number of mechanism invocations is $T$.
This is not true of record-level DP-SGD, where every local step is a separate
release.

\subsection{Adversary}

We consider a passive observer who follows the protocol but examines what it can
see. It knows the region structure and the noise scales $\{\sigma_j\}$, and it
observes the regional sums $Y_r^{(t)}$ for all rounds. It does not observe
individual silo updates, which is the property the aggregation tier provides.

The observer's goal is to infer, for a target silo $i$, the concentration
\begin{equation}
p_i \;=\; \frac{|\{x \in D_i : x \text{ belongs to the sensitive class}\}|}{|D_i|},
\end{equation}
quantised to $R$ levels so that $H(p_i)$ is finite. This target captures
information that organisations are often obliged to protect, such as the
prevalence of a condition at a particular clinical site, and it is a property of
the dataset rather than of any single record.

\section{Privacy Analysis}
\label{sec:theory}

\subsection{Lateral leakage}

Some information about $p_i$ is available to the observer without seeing any
model updates at all, because silos in the same organisational group tend to have
correlated distributions. We write this floor as
\begin{equation}
\ell_i \;=\; I\big(p_i ; D_{-i}\big),
\end{equation}
evaluated under the deployment's stated generative model, where
$D_{-i}$ denotes the datasets of all other silos. No amount of noise added to
silo $i$'s updates can reduce $\ell_i$, since it does not flow through the
mechanism at all.

It is important that $\ell_i$ be defined relative to a stated prior rather than as
a supremum over all priors consistent with the structure. Under an unconstrained
supremum one may always choose a prior making $p_i$ a deterministic function of a
neighbour's data, in which case $\ell_i = H(p_i)$ for every silo with at least one
neighbour, the quantity is constant across silos, and any allocation derived from
it degenerates to the uniform one.

Under a hierarchical model in which a group-level parameter $\Phi_g$ is drawn
first and each member's concentration is drawn conditionally on it, $\ell_i$ is
computable and depends only on the number of other members of $i$'s group. It
increases with that number and saturates at $I(p_i ; \Phi_g)$, which is strictly
below $H(p_i)$. Saturation is rapid: in our calibration the first sibling
contributes roughly two thirds of the eventual total, and the difference between a
group of twenty and a group of fifty is negligible. Consequently $\ell_i$ varies
little across silos in practice, and it is not the quantity that drives the
allocation.

\subsection{The per-client bound}

\begin{theorem}[Per-client information bound]
\label{thm:bound}
Let $\hat{p}_i = f(Y)$ be any estimate formed by a deterministic function of the
observations $Y = \{Y_r^{(t)}\}$. Then
\begin{equation}
I\big(p_i ; \hat{p}_i \,\big|\, \{\sigma_j\}\big)
\;\le\;
\min\Big\{\, H(p_i), \;\; T \cdot m_i(\sigma) + \ell_i \,\Big\},
\label{eq:bound}
\end{equation}
where the mechanism term is
\begin{equation}
m_i(\sigma) \;=\; \frac{2\,w_i^2}{\sum_{j \in A(i)} w_j^2 \sigma_j^2}.
\label{eq:mterm}
\end{equation}
\end{theorem}

The proof is given in the appendix. Its structure is as follows. The data
processing inequality reduces the estimate to the observations. A chain rule
separates the observations into what they reveal beyond $D_{-i}$ and what
$D_{-i}$ reveals on its own, the latter being $\ell_i$ by definition. Conditioning
on the observations from earlier rounds then isolates silo $i$'s contribution
within each round, since given the past every other silo's update is a
deterministic function of the previous global model and its own data, perturbed by
independent noise. The per-round contribution is bounded by the largest
Kullback-Leibler divergence the mechanism can produce between two silo-adjacent
inputs, and these bounds add across rounds.

Three properties of \eqref{eq:bound} are worth drawing out.

The mechanism term depends on the total noise in the region, not on silo $i$'s
noise alone. This is the formal statement of the intuition that a participant is
concealed by its neighbours. For $m$ equally weighted silos sharing a region, all
using the same scale $\sigma$, the term reduces to $2/(m \sigma^2)$, so the noise
each silo must contribute to reach a given level falls as $1/\sqrt{m}$.

The bound enters through the region, so the region structure is something the
mechanism can act on. This is what makes the allocation problem in
Section~\ref{sec:allocation} non-trivial.

The cap at $H(p_i)$ makes it possible to say whether the guarantee is
informative. The bound conveys something only when
$T \cdot m_i(\sigma) < H(p_i) - \ell_i$. We refer to the right-hand side as the
margin, since it is the amount of information about $p_i$ that is not already
determined by the deployment structure and that the mechanism is therefore
responsible for protecting. Reporting a bound of several nats against a target
whose entropy is smaller conveys nothing, and we check this condition explicitly
at every operating point we evaluate.

\section{Optimal Allocation}
\label{sec:allocation}

\subsection{The problem}

Let $S_r = \sum_{j \in r} w_j^2 \sigma_j^2$ denote the noise variance that region
$r$ contributes to the aggregate. The utility cost of privacy is the total noise
variance entering the global model,
\begin{equation}
U \;=\; \sum_i w_i^2 \sigma_i^2 \;=\; \sum_r S_r ,
\label{eq:budget}
\end{equation}
and not $\sum_i \sigma_i^2$. The distinction matters under weighted averaging,
where a silo's noise enters the aggregate scaled by its weight. Using the
unweighted sum both misstates the cost and makes the optimisation degenerate,
because it becomes cheapest to concentrate all of a region's noise on its
highest-weight member.

Writing $W_r = \max_{i \in r} w_i^2$ for the weight of the most exposed silo in
region $r$, the worst-case bound within the region is
$a W_r / S_r + \ell_r$ with $a = 2T$. The allocation problem is therefore
\begin{equation}
\min_{\{S_r > 0\}} \; \max_r \left[ \frac{a W_r}{S_r} + \ell_r \right]
\quad \text{subject to} \quad \sum_r S_r \le U .
\label{eq:opt}
\end{equation}

\subsection{Solution}

\begin{theorem}[Min-max optimal allocation]
\label{thm:allocation}
For any budget $U > 0$, the solution of \eqref{eq:opt} is
\begin{equation}
S_r^{\star} \;=\; \frac{a W_r}{K^{\star} - \ell_r},
\label{eq:alloc}
\end{equation}
where $K^{\star}$ is the unique root of
$\sum_r a W_r / (K - \ell_r) = U$ on $(\max_r \ell_r, \infty)$. The achieved
worst-case bound equals $K^{\star}$ and is equalised across regions.
\end{theorem}

The objective is not smooth, but introducing a slack variable for the maximum
converts \eqref{eq:opt} into a convex program with a linear objective. Slater's
condition holds, so the Karush-Kuhn-Tucker conditions are necessary and
sufficient, and they give \eqref{eq:alloc} directly. Uniqueness of $K^{\star}$
follows because the left-hand side of the budget equation is continuous and
strictly decreasing on the relevant interval, tends to infinity at its left
endpoint, and tends to zero at infinity. The full argument is in the appendix.

\subsection{How much is gained}

Uniform allocation applies one scale $\sigma$ to every silo, so that
$S_r = \sigma^2 V_r$ with $V_r = \sum_{j \in r} w_j^2$. To meet a target $K$ it
must satisfy the most demanding region, which fixes $\sigma^2$ and hence the total
budget. Comparing the two budgets required to reach the same target gives the
following.

\begin{corollary}[Budget saved at equal privacy]
\label{cor:delta}
Define the exposure ratio of region $r$ as
\begin{equation}
\rho_r \;=\; \frac{W_r}{V_r} \;=\; \frac{\max_{i \in r} w_i^2}{\sum_{i \in r} w_i^2},
\end{equation}
and let $\langle \rho \rangle_V = \sum_r \rho_r V_r / \sum_r V_r$. If the lateral
floors $\ell_r$ are equal across regions, the fraction of the noise budget saved
by \eqref{eq:alloc} relative to uniform allocation at the same target is
\begin{equation}
\delta \;=\; 1 - \frac{\langle \rho \rangle_V}{\rho_{\max}} .
\label{eq:delta}
\end{equation}
\end{corollary}

The quantity $\rho_r$ is the inverse participation ratio of the weights within
region $r$. For $m$ equally weighted silos it equals $1/m$, so it measures the
region's effective size, and it accounts correctly for the case where a region is
dominated by one large participant. Three consequences follow.

First, $\delta$ is determined entirely by the region-size profile. It does not
depend on the dataset, the model, the number of rounds, or the privacy target.
A deployer can therefore compute it from the deployment description before
training anything and decide whether the method is worth adopting.

Second, $\delta = 0$ exactly when $\rho_r$ is the same for every region, that is
when all regions have the same effective size. The method then provably does
nothing, and we treat this as an applicability test to be stated up front rather
than a limitation to be discovered later.

Third, flat federated learning, in which each silo forms its own region and its
update is observed individually, has $\rho_r = 1$ for every silo and therefore
$\delta = 0$ regardless of how unequal the silo sizes are. Without an aggregation
tier there is nothing to reallocate. The scope of this paper is consequently
hierarchical deployments.

\subsection{Exposure in real deployments}

Because $\delta$ is a property of the deployment rather than of an experiment, it
can be measured directly for real federations. Table~\ref{tab:realdelta} reports
it for several. Cellular deployments sit high on the scale because base station
loads are long tailed, with dense cells serving many devices and sparse cells very
few. Consortia organised by country behave similarly. A federation whose sites are
grouped into equally sized regions sits at zero, as does flat cross-silo training.

\begin{table}[t]
\caption{Exposure dispersion $\delta$ for real deployment structures. $\delta$ upper-bounds the fraction of the noise budget that allocation can save.}
\label{tab:realdelta}
\centering
\small
\begin{tabular}{@{}lcc@{}}
\toprule
Deployment structure & Regions & $\delta$ \\
\midrule
Cellular edge, long-tailed devices per cell & 60 & $82.8$--$87.6\%$ \\
Consortium, sites per country $30/10/5/3/1/1$ & 6 & $88.0\%$ \\
Clinical sites grouped by hospital $3/1/1/1$ & 4 & $14.4\%$ \\
Equally sized regions & any & $0\%$ \\
Flat federation, no aggregation tier & --- & $0\%$ \\
\bottomrule
\end{tabular}
\end{table}

\section{Evaluation}
\label{sec:eval}

\subsection{Design}

Corollary~\ref{cor:delta} makes a prediction that an experiment can contradict,
so we designed the evaluation around that possibility rather than around a
favourable configuration. We sweep the region-size profile across the range of
$\delta$, including two deployments with $\delta = 0$ for which the theory
predicts no gain whatsoever. We state three predictions in advance.

The first is that the budget saved at a matched privacy target equals $\delta$.
The second is that the accuracy gain at matched privacy increases with $\delta$
and is zero when $\delta = 0$. The third is that an allocation with the same
dispersion of noise but assigned to the wrong regions does not reproduce the gain,
which rules out the possibility that any non-uniform allocation would do as well.

All three arms hold every silo at the same worst-case bound $K$, so the comparison
is at matched per-client privacy and the arms differ only in the total noise they
require. We report the privacy parameters achieved in every run. Client $i$ is
concealed by the total noise in its region, so its effective noise multiplier is
$\sqrt{S_r}/w_i$, and we compute $\varepsilon$ from the least protected client.

\subsection{Setup}

We use $n = 96$ silos, $T = 10$ rounds, $\delta_{\mathrm{DP}} = 10^{-5}$, clipping
at $C = 1$, and three seeds per configuration. The privacy target is $K = 0.88$
nats, which lies below the entropy of the quantised target and is therefore
informative in the sense of Section~\ref{sec:theory}. The resulting privacy
parameter is $\varepsilon = 0.99$ in every run reported here.

Both tasks are federated fine-tuning of a frozen pretrained backbone. For images
we take CIFAR-10 grouped into two superclasses, embed it with an
ImageNet-pretrained ResNet-18, reduce to 32 dimensions, and train a linear head.
For text we take AG News reduced to two classes, embed it with a frozen
MiniLM sentence encoder, and apply the same reduction and head. The head
dimension is identical across the two tasks, which holds the signal-to-noise
regime fixed and makes the results comparable.

The choice of fine-tuning rather than training from scratch is deliberate and is
discussed in Section~\ref{sec:discussion}.

\subsection{Results}

Table~\ref{tab:main} reports the outcome. The budget saved matches $\delta$ to
three decimal places in every configuration, confirming the first prediction and
Corollary~\ref{cor:delta}.

The accuracy gain increases with $\delta$, reaching $13.75$ points on images and
$13.55$ points on text in the most heterogeneous deployment, which confirms the
second prediction. In the two balanced deployments the gain is exactly zero, to
the precision of the measurement, on both tasks. This is the outcome the theory
requires, and we regard it as the most informative row in the table: a method that
did not act on the quantity we claim it acts on would be unlikely to produce
exactly zero here.

The misallocation control confirms the third prediction. Given the same dispersion
of noise but assigned to the wrong regions, and rescaled so that it still meets
the same target, it performs worse than uniform allocation by between two and
three points in every non-trivial configuration. Non-uniformity alone is therefore
not sufficient, and the gain depends on allocating according to the exposure
profile.

Agreement between the two modalities is close, at $4.33$ against $4.02$ points and
$13.75$ against $13.55$ points. We had expected the magnitudes to differ, since
the gain depends on the local slope of accuracy against noise, which is a property
of the task. The close agreement is likely a consequence of having matched the
head dimension and privacy target across the two tasks, and we would not expect it
to persist if either were varied.

\begin{table*}[t]
\caption{Accuracy at matched per-client privacy ($\varepsilon = 0.99$, $n = 96$, $T = 10$, three seeds). Gains are percentage points relative to uniform allocation. Non-private references are $0.970$ for CIFAR-10 and $0.930$ for AG News.}
\label{tab:main}
\centering
\small
\begin{tabular}{@{}lcc cccc cccc@{}}
\toprule
& & & \multicolumn{4}{c}{CIFAR-10 (frozen ResNet-18)} & \multicolumn{4}{c}{AG News (frozen MiniLM)} \\
\cmidrule(lr){4-7}\cmidrule(lr){8-11}
Region profile & $\delta$ & Budget saved
& \Fulcrum & Uniform & Gain & Misalloc.
& \Fulcrum & Uniform & Gain & Misalloc. \\
\midrule
Balanced, 4 regions (control) & $0.000$ & $0.0\%$
& $0.868$ & $0.868$ & $+0.00$ & $+0.00$
& $0.767$ & $0.767$ & $+0.00$ & $+0.00$ \\
Balanced, 6 regions (control) & $0.000$ & $0.0\%$
& $0.832$ & $0.832$ & $+0.00$ & $+0.00$
& $0.734$ & $0.734$ & $+0.00$ & $+0.00$ \\
Mild, $6/6/6/3/3$ & $0.375$ & $37.5\%$
& $0.851$ & $0.808$ & $+4.33$ & $-3.12$
& $0.748$ & $0.708$ & $+4.02$ & $-2.55$ \\
Severe, $15/5/2/1/1$ & $0.792$ & $79.2\%$
& $0.845$ & $0.708$ & $\mathbf{+13.75}$ & $-2.33$
& $0.759$ & $0.623$ & $\mathbf{+13.55}$ & $-3.25$ \\
\bottomrule
\end{tabular}
\end{table*}

\section{Discussion}
\label{sec:discussion}

\subsection{Why fine-tuning}

The noise added to the aggregate has norm proportional to
$C \sigma \sqrt{d \sum_i w_i^2}$, where $d$ is the number of trainable
parameters, while the aggregate signal has norm proportional to $C \sqrt{a}$ for
an average alignment $a$ among the clipped updates. The ratio therefore behaves as
$\sqrt{a n / d} \, / \, \sigma$. Parameter count enters directly, and at
silo-level granularity it is the binding constraint.

The practical consequence is sharp. A convolutional network trained from scratch
in our setting has $d$ on the order of $5 \times 10^5$ against $n = 96$ silos, and
at any noise level for which the bound is informative the model does not learn at
all. Fine-tuning a small head on a frozen backbone reduces $d$ by four orders of
magnitude and reaches full non-private accuracy at the same privacy target. We
therefore present silo-level differential privacy as suitable for federated
fine-tuning rather than for federated training from scratch, and we regard this as
a deployment guideline rather than a limitation of the allocation result, which
holds regardless.

\subsection{Scope and assumptions}

The allocation is worth applying only when regions differ in effective size, and
$\delta$ quantifies this before deployment. Flat federations gain nothing.

The analysis assumes disjoint regions, static membership over training, and
disjoint silo datasets. Overlapping regions, in which a silo contributes to
several aggregates, would require the observer's view to be modelled jointly
across regions and is not covered.

The bound is an upper bound and is loose in at least two places. The chain rule
step discards a non-negative term, and the per-round divergence is a worst case
over adjacent inputs. Tightening either would reduce the noise required at a given
target, but would not change the allocation, since the allocation depends on the
relative exposure of regions rather than on the absolute level of the bound.

We assume a passive observer. An active adversary that manipulates the protocol,
or a malicious aggregator that misreports sums, is outside the model. Because each
silo adds its own noise before transmitting, the guarantee does not require
trusting the aggregator to add noise, but it does assume the aggregator reports
the sum it computed.

\subsection{Relation to trusted-aggregator schemes}

As noted in Section~\ref{sec:related}, a trusted super-node that injects a single
noise draw on behalf of its region requires per-client noise scaling as $1/m$
rather than $1/\sqrt{m}$. A deployment able to trust its aggregators can therefore
reach the same guarantee with less noise. The allocation problem we solve arises
in both settings, and Theorem~\ref{thm:allocation} applies to either once the
mechanism term is written accordingly, but the value of $\delta$ in
Corollary~\ref{cor:delta} is derived for the untrusted case treated here.

\section{Conclusion}
\label{sec:conclusion}

In hierarchical federated learning a participant is concealed by the region it is
aggregated with, and regions in real deployments differ substantially in effective
size. Uniform noise must be calibrated for the most exposed region and therefore
over-provisions the rest. We derived a per-client information bound whose
adjacency matches the property being protected, solved the resulting allocation
problem in closed form, and showed that the saving relative to uniform allocation
is exactly $\delta = 1 - \langle \rho \rangle_V / \rho_{\max}$, a quantity
computable from the deployment before training and zero precisely when regions are
equally sized. Experiments on two modalities confirmed the prediction at
$\varepsilon = 0.99$, including two controls for which the theory predicts no
gain and which delivered none.

Two directions follow naturally. An allocation combining structural exposure with
declared per-participant preferences would unify this work with personalised
differential privacy. Extending the analysis to overlapping regions would cover
deployments in which a silo contributes to more than one aggregate, which the
present treatment excludes.

\section*{Data and Code Availability}

Source code and the analysis artefacts underlying every table are available at
\url{https://doi.org/10.5281/zenodo.20507155}.

\begin{thebibliography}{99}
\bibitem{abadi2016dpsgd} M. Abadi \textit{et al.}, ``Deep learning with differential privacy,'' in \textit{ACM CCS}, 2016, pp. 308--318.
\bibitem{mcmahan2018dprnn} H. B. McMahan, D. Ramage \textit{et al.}, ``Learning differentially private recurrent language models,'' in \textit{ICLR}, 2018.
\bibitem{mironov2017rdp} I. Mironov, ``R\'enyi differential privacy,'' in \textit{IEEE CSF}, 2017, pp. 263--275.
\bibitem{jorgensen2015personalized} Z. Jorgensen, T. Yu, and G. Cormode, ``Conservative or liberal? Personalized differential privacy,'' in \textit{IEEE ICDE}, 2015, pp. 1023--1034.
\bibitem{chandrasekaran2022hfl} V. Chandrasekaran, S. Banerjee, D. Perino, and N. Kourtellis, ``Hierarchical federated learning with privacy,'' \textit{arXiv:2206.05209}, 2022.
\bibitem{guo2025fedcdp} P. Guo, C. Bai, M. Zhang, and P. O. Dusadeerungsikul, ``Clustered federated learning with heterogeneous differential privacy on non-IID data,'' \textit{Computer Communications}, vol. 244, p. 108339, 2025.
\bibitem{yin2026hdpfedcd} C. Yin, K. He, and J. Shi, ``HDP-FedCD: Data-quality-driven hierarchical federated learning for optimizing privacy protection in non-IID data,'' \textit{Future Generation Computer Systems}, vol. 176, p. 108140, 2026.
\bibitem{bonawitz2017secagg} K. Bonawitz \textit{et al.}, ``Practical secure aggregation for privacy-preserving machine learning,'' in \textit{ACM CCS}, 2017, pp. 1175--1191.
\bibitem{cyffers2022amplification} E. Cyffers and A. Bellet, ``Privacy amplification by decentralization,'' in \textit{AISTATS}, 2022, pp. 5334--5353.
\bibitem{cyffers2022muffliato} E. Cyffers, M. Even, A. Bellet, and L. Massouli\'e, ``Muffliato: Peer-to-peer privacy amplification for decentralized optimization and averaging,'' in \textit{NeurIPS}, 2022.
\bibitem{krizhevsky2009cifar} A. Krizhevsky, ``Learning multiple layers of features from tiny images,'' Tech. Rep., 2009.
\bibitem{zhang2015agnews} X. Zhang, J. Zhao, and Y. LeCun, ``Character-level convolutional networks for text classification,'' in \textit{NeurIPS}, 2015.
\end{thebibliography}

\appendices

\section{Proof of Theorem~\ref{thm:bound}}
\label{app:bound}

We first record two standard facts.

\begin{lemma}[Chain rule bound]
\label{lem:chain}
For random variables $X, Y, Z$, $I(X;Y) \le I(X;Z) + I(X;Y \mid Z)$.
\end{lemma}
\begin{proof}
Expanding $I(X;Y,Z)$ in two ways gives
$I(X;Z) + I(X;Y \mid Z) = I(X;Y) + I(X;Z \mid Y)$,
and $I(X;Z \mid Y) \ge 0$.
\end{proof}

\begin{lemma}[Maximal divergence bound]
\label{lem:maxkl}
If a channel $P_{Y \mid X}$ satisfies
$\KL(P_{Y \mid x} \,\|\, P_{Y \mid x'}) \le \kappa$ for all $x, x'$, then
$I(X;Y) \le \kappa$.
\end{lemma}
\begin{proof}
$I(X;Y) = \mathbb{E}_X \KL(P_{Y \mid X} \| P_Y)$ with
$P_Y = \mathbb{E}_{X'} P_{Y \mid X'}$. By convexity of the divergence in its
second argument, $\KL(P_{Y \mid x} \| P_Y) \le \mathbb{E}_{X'} \KL(P_{Y \mid x} \| P_{Y \mid X'}) \le \kappa$.
\end{proof}

\begin{proof}[Proof of Theorem~\ref{thm:bound}]
Since $\hat{p}_i$ is a deterministic function of $Y$, the data processing
inequality gives $I(p_i ; \hat{p}_i) \le I(p_i ; Y)$. Applying
Lemma~\ref{lem:chain} with $Z = D_{-i}$,
\begin{equation}
I(p_i ; Y) \;\le\; I(p_i ; D_{-i}) + I(p_i ; Y \mid D_{-i})
\;=\; \ell_i + I(p_i ; Y \mid D_{-i}).
\end{equation}
Because $p_i$ is a function of $D_i$, a further application of the data
processing inequality gives
$I(p_i ; Y \mid D_{-i}) \le I(D_i ; Y \mid D_{-i})$.

Expanding over rounds by the chain rule,
\begin{equation}
I(D_i ; Y \mid D_{-i}) \;=\; \sum_{t=1}^{T} I\big(D_i ; Y^{(t)} \,\big|\, Y^{(<t)}, D_{-i}\big).
\end{equation}
Fix a round $t$ and condition on $Y^{(<t)}$ and $D_{-i}$. The global model
$\theta^{(t-1)}$ is a deterministic function of $Y^{(<t)}$ through the aggregation
rule. Hence for every $j \ne i$ the update $\Delta_j^{(t)}$ is a deterministic
function of $\theta^{(t-1)}$ and $D_j$, and the noise $\xi_j^{(t)}$ is drawn
independently. Writing the region sum containing $i$ as
\begin{equation}
Y_{A(i)}^{(t)} \;=\; w_i \big(\mathrm{clip}(\Delta_i^{(t)}, C) + \xi_i^{(t)}\big) \;+\; R,
\end{equation}
the remainder $R$ is a function of $Y^{(<t)}$, $D_{-i}$ and noise independent of
$D_i$, and is therefore conditionally independent of $D_i$ given the conditioning
variables. Region sums not containing $i$ are likewise conditionally independent
of $D_i$. We note that no claim of unconditional independence between $D_i$ and
the other silos' updates is made or needed; conditioning on the past is what makes
the step valid, since across rounds the updates of other silos do depend on $D_i$
through the shared global model.

The conditional channel from $D_i$ to $Y_{A(i)}^{(t)}$ is therefore Gaussian with
covariance $\big(\sum_{j \in A(i)} w_j^2 \sigma_j^2\big) C^2 I$ and with mean
shifted by $w_i \,\mathrm{clip}(\Delta_i^{(t)}, C)$. Replacing $D_i$ by any other
dataset moves the clipped update by at most $2C$ in norm, and hence moves the mean
by at most $2 w_i C$. The divergence between two such Gaussians is
\begin{equation}
\frac{(2 w_i C)^2}{2 \big(\sum_{j \in A(i)} w_j^2 \sigma_j^2\big) C^2}
\;=\; \frac{2 w_i^2}{\sum_{j \in A(i)} w_j^2 \sigma_j^2} \;=\; m_i(\sigma),
\end{equation}
so Lemma~\ref{lem:maxkl} bounds each round's term by $m_i(\sigma)$. Summing over
$T$ rounds and adding $\ell_i$ gives the second entry of the minimum. The first
entry holds because $I(p_i ; \hat{p}_i) \le H(p_i)$ for any $\hat{p}_i$.

Both the left-hand side and the divergence computation refer to replacement of the
entire dataset $D_i$, so the adjacency notion is the same throughout and no group
privacy argument is required.
\end{proof}

\section{Proof of Theorem~\ref{thm:allocation}}
\label{app:alloc}

Introduce a slack variable $K$ and rewrite \eqref{eq:opt} as
\begin{equation}
\min_{K, \{S_r > 0\}} K
\quad \text{s.t.} \quad
\frac{a W_r}{S_r} + \ell_r \le K \;\; \forall r,
\qquad \sum_r S_r \le U .
\label{eq:convex}
\end{equation}
The objective is linear and each constraint is convex, since $S \mapsto aW_r/S$ is
convex and decreasing on $S > 0$. Setting $S_r = U / (2 |\mathcal{R}|)$ for all
$r$, where $|\mathcal{R}|$ is the number of regions, is strictly feasible for $K$
large enough, so Slater's condition holds and the Karush-Kuhn-Tucker conditions
are necessary and sufficient.

The Lagrangian is
\begin{equation}
\mathcal{L} = K + \sum_r \mu_r \left( \frac{a W_r}{S_r} + \ell_r - K \right)
+ \lambda \left( \sum_r S_r - U \right),
\end{equation}
with $\mu_r \ge 0$ and $\lambda \ge 0$. Stationarity in $K$ gives
$\sum_r \mu_r = 1$. Stationarity in $S_r$ gives
$-\mu_r a W_r / S_r^2 + \lambda = 0$, hence
$S_r = \sqrt{\mu_r a W_r / \lambda}$.

Every per-region constraint is active at the optimum. Suppose otherwise for some
$r$; complementary slackness then forces $\mu_r = 0$, and stationarity in $S_r$
forces $\lambda = 0$, which in turn forces $\mu_{r'} = 0$ for all $r'$ and
contradicts $\sum_r \mu_r = 1$. Activity of every constraint gives
$a W_r / S_r + \ell_r = K^{\star}$ for all $r$, which rearranges to
\eqref{eq:alloc}.

For uniqueness, define $g(K) = \sum_r a W_r / (K - \ell_r)$ on
$(\max_r \ell_r, \infty)$. It is continuous and strictly decreasing there, tends
to $+\infty$ as $K$ approaches the left endpoint, and tends to $0$ as
$K \to \infty$. By the intermediate value theorem there is exactly one
$K^{\star}$ with $g(K^{\star}) = U$.

\section{Proof of Corollary~\ref{cor:delta}}
\label{app:delta}

Assume $\ell_r = \ell$ for all $r$ and write $c = a / (K - \ell)$ for a target
$K > \ell$. By \eqref{eq:alloc} the optimal allocation requires
$U_{\mathrm{opt}} = c \sum_r W_r$.

Uniform allocation uses a single scale $\sigma$, giving $S_r = \sigma^2 V_r$. The
constraint for region $r$ reads $a W_r / (\sigma^2 V_r) + \ell \le K$, that is
$\sigma^2 \ge c \, W_r / V_r = c \rho_r$. Meeting every region therefore requires
$\sigma^2 = c \, \rho_{\max}$, and the resulting budget is
$U_{\mathrm{uni}} = c \, \rho_{\max} \sum_r V_r$.

Using $W_r = \rho_r V_r$,
\begin{equation}
\frac{U_{\mathrm{opt}}}{U_{\mathrm{uni}}}
= \frac{\sum_r \rho_r V_r}{\rho_{\max} \sum_r V_r}
= \frac{\langle \rho \rangle_V}{\rho_{\max}},
\end{equation}
and $\delta = 1 - U_{\mathrm{opt}} / U_{\mathrm{uni}}$ gives \eqref{eq:delta}.
Since $\rho_r \le \rho_{\max}$ for every $r$, the ratio is at most one and
$\delta \ge 0$, with equality if and only if $\rho_r$ is constant across regions.

\end{document}
